% Generated draft: requires mathematical and source review.
\documentclass[11pt]{article}
\usepackage[margin=25mm]{geometry}
\usepackage{fontspec}
\setmainfont{Latin Modern Roman}
\usepackage{amsmath,amssymb,mathtools,unicode-math}
\setmathfont{Latin Modern Math}
\usepackage{enumitem,needspace,xurl,hyperref}
\setlength{\parindent}{0pt}
\setlength{\parskip}{5pt}
\setlength{\emergencystretch}{3em}
\newcommand{\N}{\mathbb{N}}
\newcommand{\Z}{\mathbb{Z}}
\newcommand{\Q}{\mathbb{Q}}
\newcommand{\R}{\mathbb{R}}
\newcommand{\C}{\mathbb{C}}
\newcommand{\F}{\mathbb{F}}
\newcommand{\A}{\mathbb{A}}
\newcommand{\PP}{\mathbb{P}}
\newcommand{\E}{\mathbb{E}}
\newcommand{\eps}{\varepsilon}
\newcommand{\dd}{\,\mathrm{d}}
\begin{document}
\section*{UnsolvedMath ID: 30001920}
\textbf{Source problem code:} catalogue entry
\subsection{Definitions and mathematical statement}
Write \((\widetilde X_\sigma,w)\) and \((\widetilde X_{\sigma^\vee},w^\vee)\) for a mirror pair in the class intended by the supplied source record. Here \(\sigma\) and \(\sigma^\vee\) label the two constructions, and \(w,w^\vee\) are their accompanying data. The record does not specify these constructions, the meaning of mirror pair, or the domains and targets of \(w,w^\vee\). Let \(n\) denote the parameter in the supplied formula; its precise dimensional meaning is not supplied. The notation \(h^{p,q}(\widetilde X_\sigma,w)\) denotes the intended relative Hodge number, whose underlying cohomology theory and indexing conventions are not defined in the record.

\textbf{Mirror Hodge-number conjecture, as supplied.} For every such mirror pair for which both \(\widetilde X_\sigma\) and \(\widetilde X_{\sigma^\vee}\) are smooth, the proposed identity is
\[
 h^{p,q}(\widetilde X_\sigma,w)
 =h^{n-p,q}(\widetilde X_{\sigma^\vee},w^\vee)
\]
for every bidegree \((p,q)\) in the intended indexing range. In the orbifold case, the source requests the same exchange with orbifold Hodge numbers in place of these Hodge numbers.

The intended class of mirror pairs, any additional hypotheses, the value and meaning of \(n\), the allowed bidegrees, the precise relative Hodge-number construction, and its orbifold analogue remain unspecified. In particular, no relative cohomology model or orbifold grading convention is imposed here. These omissions prevent a fully specified mathematical conjecture; the displayed identity preserves the supplied formulation pending human mathematical review.
\subsection{Short English statement}
The supplied conjecture predicts that relative Hodge numbers of a smooth mirror pair agree after replacing the first Hodge index \(p\) by \(n-p\). It also proposes an orbifold version using orbifold Hodge numbers. The geometric hypotheses and the definitions of these invariants require clarification.
\subsection{Sources}
\begin{itemize}
\item Supplied local UnsolvedMath registry record, permanent ID \(30001920\), title \emph{Mirror Symmetry for Relative Hodge Numbers}, source problem code \texttt{catalogue entry}, in the Oberwolfach Reports Open Problems collection. The record gives proposed year 2011 and names no individual proposer.
\item Supplied report citation: \emph{Complex Algebraic Geometry (2012)}, Oberwolfach Reports, pp.~2669--2733, DOI \texttt{10.4171/owr/2011/47}. Original OWR report: \url{https://doi.org/10.4171/owr/2011/47}. The supplied title year is 2012, while the proposed year and DOI identifier refer to 2011; their relationship has not been independently checked.
\item Supplied literature reference, labeled ``published mirror-symmetry theorem'': \url{https://doi.org/10.1016/j.aim.2016.03.035}. The supplied assessment attributes to Gross--Katzarkov--Ruddat a proof of the expected Hodge-number exchange for their toric hypersurface construction, while describing the broader orbifold and general mirror formulation as only partially established. The record does not provide the precise theorem or its hypotheses.
\item Literature assessment dated exactly \texttt{2026-08-21}: the supplied status is \texttt{partially\_solved}. This is supplied dated curation, not an independently verified literature assessment. Both supplied literature-source links are retained above.
\item The registry labels the statement \texttt{corrected\_verified} and reports that its mathematical content was checked against the official Oberwolfach report, with notation, context, and question form normalized without changing the source-level claim. This is a supplied verification claim; no source visit or independent verification is asserted here.
\end{itemize}
\end{document}
